Savstarpēji savienoti \Gls{lif} neironi digitālā arhitektūrā ir plaši izmantota izvēle liela mēroga \Gls{snn} ar efektīvu
resursu patēriņu, tomēr sarežģītākas impulsu ģenerēšanas dinamikas izpētei sistēmai nepieciešams augstāks bioloģiskās
līdzības līmenis. Lai precīzāk atdarinātu neirona uzvedību, parasti izvēlas reālistiskākus modeļus, kas spēj radīt dažādus
pulsēšanas režīmus, piemēram, Ižikeviča modeli \cite{Izhikevich2003}, taču tas palielina skaitļošanas sarežģītību.

Ižikeviča neirona modelī \cite{Izhikevich2003} Hodžkina–Hakslija neirona modelis \cite{hodgkin1952} reducēts līdz šādai
parasto diferenciālvienādojumu sistēmai:
\begin{equation}
\begin{cases}
  \label{eq:izhi}
  v' = 0{,}04v^{2} + 5v + 140 - u + I \\
  u' = a(bv - u),
\end{cases}
\end{equation}
kur \(v\) apzīmē membrānas potenciālu, \(u\) ir membrānas atjaunošanās mainīgais, \(a\) nosaka atjaunošanās mainīgā laika
mērogu un \(b\) raksturo jutību pret zemsliekšņa svārstībām.

Kad impulss sasniedz potenciāla slieksni, sistēma tiek atiestatīta uz pēcimpulsa stāvokli:
\begin{equation}
  \label{eq:postspike}
  v \geq V_{th} \text{, tad } \{v = c, u = u + d\},
\end{equation}
kur parametri \(c\) un \(d\) raksturo membrānas potenciāla un atjaunošanās mainīgā pēcimpulsa atiestatīšanas vērtības.
Pretstatā bioloģiski nereālistiski vienkāršajam \Gls{lif} modelim šis modelis spēj atveidot smadzeņu garozas neironiem raksturīgo atsevišķu
impulsu un impulsu kopu dinamiku.


